\documentclass[12pt,a4paper]{report}

\usepackage[utf8]{inputenc}
\usepackage[T1]{fontenc}
\usepackage[turkish,english]{babel}
\usepackage{graphicx}
\usepackage{float}
\usepackage{amsmath}
\usepackage{array}
\usepackage{booktabs}
\usepackage{longtable}
\usepackage{url}
\usepackage{hyperref}
\usepackage{setspace}
\usepackage{geometry}

\geometry{left=3cm,right=2.5cm,top=3cm,bottom=3cm}
\onehalfspacing

\begin{document}

%-------------------- KAPAK --------------------
\begin{titlepage}
    \thispagestyle{empty}
    \begin{center}
        \textbf{T.C.}\\[0.3cm]
        \textbf{İSTANBUL AREL ÜNİVERSİTESİ}\\[0.2cm]
        \textbf{MÜHENDİSLİK FAKÜLTESİ}\\[0.2cm]
        \textbf{BİLGİSAYAR MÜHENDİSLİĞİ BÖLÜMÜ}\\[0.8cm]

        \includegraphics[width=2.4cm,trim=2cm 2cm 2cm 2cm,clip]{logo.png}\\[0.8cm]

        \textbf{\Large Türkiye Deprem Verileri Kullanılarak Riskli Bölge\\
        ve Olay Deseni Analizi için Karar Destek Sistemi}\\[1.5cm]

        \textbf{Arman Yalçın}\\
        \textbf{220309012}\\[1.2cm]

        \textbf{Danışman}\\
        \textbf{Dr. Sibel BİRTANE AKAR}\\[1.5cm]

        \textbf{ARA RAPOR}\\[1.5cm]

        İstanbul, Haziran 2026
    \end{center}
\end{titlepage}

%-------------------- ÖNSÖZ --------------------
\chapter*{ÖNSÖZ}
\addcontentsline{toc}{chapter}{ÖNSÖZ}

Bu ara rapor, ``Türkiye Deprem Verileri Kullanılarak Riskli Bölge ve Olay Deseni
Analizi için Karar Destek Sistemi'' başlıklı bitirme projesinin mevcut durumunu,
şimdiye kadar tamamlanan çalışmaları ve izlenecek sonraki adımları sunmak amacıyla
hazırlanmıştır. Çalışmanın bu aşamasında proje konusu netleştirilmiş, veri kaynakları
incelenmiş, veritabanı tasarımı gerçekleştirilmiş ve veri işleme sürecinin ilk çalışan
bileşenleri oluşturulmuştur.

Bu süreçte deprem verilerinin yalnızca ham olay listeleri şeklinde sunulmasının ötesine
geçerek, bölgesel yoğunluk, büyüklük dağılımı, derinlik özellikleri ve zamansal değişimlerin
birlikte değerlendirilebildiği bir karar destek yapısı geliştirmenin hem teknik hem de
akademik açıdan anlamlı olduğu görülmüştür. Hazırlanan bu ara rapor, projenin mevcut
ilerleme düzeyini ve bundan sonraki geliştirme adımlarını sistematik biçimde ortaya koymaktadır.

\vspace{1cm}
\begin{flushright}
Arman Yalçın\\
Haziran 2026
\end{flushright}

\clearpage

%-------------------- TÜRKÇE ÖZET --------------------
\chapter*{Özet}
\addcontentsline{toc}{chapter}{Özet}

Bu ara raporda, Türkiye'deki deprem kayıtlarını kullanarak riskli bölge ve olay deseni
analizi yapmayı amaçlayan bir karar destek sisteminin geliştirme sürecinde şimdiye kadar
gerçekleştirilen çalışmalar sunulmaktadır. Projenin temel amacı, AFAD ve ilerleyen aşamada
Kandilli Rasathanesi tarafından yayımlanan açık deprem verilerini kullanarak konum,
büyüklük, derinlik ve zaman boyutlarında anlamlı analizler üreten bir sistem geliştirmektir.
Proje, deprem tahmini yapmaktan ziyade geçmiş ve güncel kayıtlar üzerinden analitik
çıkarımlar sunmayı hedeflemektedir.

Ara rapor kapsamında PostgreSQL tabanlı veritabanı yapısı oluşturulmuş, deprem kayıtlarının
saklanacağı temel tablolar tasarlanmış ve AFAD verilerinin alınmasına yönelik Python tabanlı
bir veri toplama süreci geliştirilmiştir. Elde edilen veriler üzerinde tarih-saat dönüştürme,
sayısal alan temizleme ve yer bilgisinden il/ilçe ayrıştırma işlemleri gerçekleştirilmiş,
ardından veriler veritabanına aktarılmıştır. Veritabanı üzerinde şehir bazında deprem sayısı,
ortalama büyüklük, maksimum büyüklük ve örnek risk puanı hesaplamaları yapılmıştır.
Ayrıca elde edilen sonuçlar grafiksel olarak görselleştirilmiştir.

Elde edilen mevcut bulgular, projenin teknik olarak uygulanabilir olduğunu ve Türkiye'ye ait
resmi deprem verilerinin tek sistem altında analiz edilmesinin yararlı sonuçlar üretebileceğini
göstermektedir. Sonraki aşamalarda Kandilli verilerinin sisteme dahil edilmesi, risk puanlama
yaklaşımının geliştirilmesi, zaman serisi ve kümelenme analizlerinin eklenmesi ve web tabanlı
görselleştirme katmanının oluşturulması planlanmaktadır.

\clearpage

%-------------------- İNGİLİZCE ÖZET --------------------
\chapter*{İngilizce Özet}
\addcontentsline{toc}{chapter}{İngilizce Özet}

This interim report presents the progress achieved so far in the development of a decision
support system designed to analyze risky regions and earthquake event patterns using earthquake
records in Türkiye. The main objective of the project is to develop a system that produces
meaningful analyses based on location, magnitude, depth, and temporal dimensions by using
open earthquake data published by AFAD and, in later stages, Kandilli Observatory. Rather than
attempting to predict earthquakes, the project focuses on generating analytical insights from
historical and recent earthquake records.

Within the scope of this interim report, a PostgreSQL-based database structure has been created,
the core tables for storing earthquake records have been designed, and a Python-based data
collection process has been developed for retrieving AFAD data. The collected data have been
processed through date-time conversion, numerical field cleaning, and province/district extraction
from location information, and then transferred into the database. Initial database analyses have
been carried out, including city-based earthquake counts, average magnitude, maximum magnitude,
and a sample risk score calculation. In addition, the obtained results have been visualized
graphically.

The findings obtained so far indicate that the project is technically feasible and that analyzing
official earthquake data from Türkiye within a single system can provide useful outcomes.
In the next stages, Kandilli data will be integrated into the system, the risk scoring approach
will be improved, time series and clustering analyses will be added, and a web-based visualization
layer will be developed.

\clearpage

%-------------------- İÇİNDEKİLER --------------------
\tableofcontents
\clearpage

\listoffigures
\addcontentsline{toc}{chapter}{Şekil Listesi}
\clearpage

\listoftables
\addcontentsline{toc}{chapter}{Tablo Listesi}
\clearpage

%-------------------- BÖLÜMLER --------------------
\chapter{Giriş}

Türkiye, aktif fay hatları üzerinde yer alması nedeniyle yüksek deprem tehlikesi taşıyan
bir ülkedir. AFAD tarafından yayımlanan Türkiye Deprem Tehlike Haritası ve güncel deprem
bilgilendirme servisleri, ülke genelindeki sismik hareketliliğin sürekli izlendiğini
göstermektedir \cite{afad_harita, afad_son}. Benzer şekilde Kandilli Rasathanesi de son
depremler ve yıllık deprem harita-grafik çalışmaları ile kamuya açık veri sağlamaktadır
\cite{kandilli_son, kandilli_yillik}. Ancak bu veriler çoğunlukla olay listeleri veya ham
kayıtlar biçiminde sunulduğundan, kullanıcıların belirli bir bölge için geçmişten bugüne
deprem davranışını doğrudan yorumlaması zorlaşmaktadır.

Bu projenin çıkış noktası, deprem verilerini yalnızca görüntüleyen değil; bu veriler üzerinde
yoğunluk, büyüklük dağılımı, derinlik özellikleri ve zamansal değişim analizi yaparak
kullanıcıya karar desteği sunan bir yapı geliştirmektir. Böylece kullanıcılar sadece
``hangi deprem ne zaman oldu?'' sorusuna değil, aynı zamanda ``hangi bölgelerde hareketlilik
yoğunlaşıyor?'', ``hangi illerde daha yüksek büyüklüklü olaylar dikkat çekiyor?'' ve
``hangi alanlar göreli olarak daha yüksek bir risk profiline sahip?'' gibi sorulara da
yanıt bulabilecektir.

Bu çalışma kapsamında Türkiye'ye ait resmi deprem verileri kullanılarak bir karar destek
sistemi geliştirilmektedir. Sistemde deprem kayıtlarının toplanması, temizlenmesi,
veritabanında saklanması ve analiz edilmesi amaçlanmaktadır. Bu doğrultuda PostgreSQL
tabanlı bir veritabanı yapısı kurulmuş, AFAD verilerinin alınmasına yönelik Python tabanlı
veri işleme süreci oluşturulmuş ve ilk analiz çıktıları elde edilmiştir. Şehir bazında
deprem sayıları, ortalama büyüklükler ve örnek risk skorları hesaplanarak projenin temel
analiz altyapısı oluşturulmuştur.

Ara raporun amacı, projenin mevcut aşamaya kadar gerçekleştirilen çalışmalarını sistematik
biçimde sunmaktır. Bu kapsamda literatür taraması, kullanılan yöntem ve araçlar, veri yapısı,
uygulama sürecinde tamamlanan adımlar ve elde edilen ilk bulgular ele alınmaktadır. Uygulamanın
nihai hâli henüz tamamlanmamış olmakla birlikte, mevcut aşamada veri toplama, veri temizleme,
veritabanına aktarma ve ilk analiz işlemleri başarıyla gerçekleştirilmiştir.

\chapter{Literatür Taraması}

Deprem odaklı karar destek sistemleri literatürde farklı amaçlarla ele alınmaktadır.
Bazı çalışmalar afet sonrası müdahale ve kaynak planlamasına odaklanırken, bazıları
risk değerlendirme, mekânsal analiz ve çok ölçütlü karar verme süreçlerini ön plana
çıkarmaktadır. Bu durum, deprem verilerinin yalnızca saklanması veya listelenmesiyle
yetinilmeyip, karar süreçlerinde anlamlı bir girdi olarak kullanılabildiğini
göstermektedir.

Deprem risk değerlendirmesinde coğrafi bilgi sistemleri ve çok ölçütlü analizlerin
birlikte kullanıldığı çalışmalar bulunmaktadır. Bu tür çalışmalarda belirli bölgelerin
deprem açısından göreli durumunun değerlendirilebilmesi için farklı ölçütlerin bir
arada ele alınması önerilmektedir. Özellikle bölgesel yoğunluk, büyüklük dağılımı,
derinlik özellikleri ve zamansal değişim gibi göstergelerin birlikte incelenmesi,
karar vericilere daha yorumlanabilir sonuçlar sunmaktadır.

Son yıllarda deprem verilerinin büyük ölçekli veri analizi ve makine öğrenmesi ile
birlikte değerlendirildiği çalışmalar da artmıştır. Bu çalışmalar, deprem verilerinin
sadece tekil olay bazında değil, desen çıkarımı, sınıflandırma ve bölgesel eğilimlerin
yorumlanması açısından da önemli bir potansiyele sahip olduğunu ortaya koymaktadır.
Özellikle büyük deprem serilerinden sonra oluşan olay kümelenmeleri, bölgesel değişimler
ve yoğunluk farkları, veri temelli yöntemlerle daha açık biçimde analiz edilebilmektedir.

Literatürdeki çalışmalar incelendiğinde, deprem alanında geliştirilen sistemlerin
çoğunlukla çok ölçütlü değerlendirme, coğrafi bilgi sistemleri, mekânsal karar destek
yaklaşımı veya veri analitiği temelli çözümler sunduğu görülmektedir. Bu proje de
benzer biçimde karar desteği üretmeyi hedeflemektedir; ancak odağını Türkiye'de kamuya
açık resmi deprem kayıtlarının analizine ve kullanıcıya anlaşılır görsel çıktılar
sunulmasına yerleştirmektedir.

Bu proje kapsamında doğrudan yapı hasarı tahmini veya afet sonrası lojistik optimizasyonu
yerine; il, ilçe ve bölge bazında deprem yoğunluğu, büyüklük profili, derinlik özellikleri
ve zamansal aktivite değişimleri üzerinden yorumlanabilir bir risk görünümü üretmek
amaçlanmaktadır. Bu yönüyle çalışma, literatürdeki karar destek yaklaşımını Türkiye'ye
özgü açık deprem verileri üzerinde uygulanabilir bir yazılım sistemine dönüştürmeyi
hedeflemektedir.

\chapter{Yöntem ve Uygulama}

Bu bölümde proje kapsamında kullanılan yöntemsel yaklaşım, tercih edilen araçlar,
oluşturulan veri akışı, belirlenen veri alanları ve uygulanan ilk risk analizi mantığı
sunulmaktadır. Ara rapor aşamasında sistemin tüm bileşenleri tamamlanmamış olmakla
birlikte, veri toplama, veri temizleme, veritabanında saklama ve ilk analiz çıktılarının
üretilmesine yönelik temel yapı oluşturulmuştur.

\section{Kullanılan Yöntem ve Araçlar}

Projede veri işleme ve analiz süreçleri için Python programlama dili kullanılmıştır.
Python'un tercih edilme nedeni; veri toplama, dönüştürme, analiz ve görselleştirme
işlemlerinin aynı ekosistem içerisinde yürütülebilmesidir. Veri temizleme ve tablo
bazlı işlemler için \texttt{pandas}, veritabanı bağlantısı için \texttt{psycopg2},
grafiksel çıktıların üretilmesi için ise \texttt{matplotlib} kullanılmıştır.

Verilerin kalıcı olarak tutulması için PostgreSQL tabanlı bir ilişkisel veritabanı
tercih edilmiştir. Bu tercihin nedeni, deprem olaylarının tarih, konum, derinlik ve
büyüklük gibi yapısal alanlarla düzenli biçimde saklanabilmesi ve analiz sorgularının
verimli şekilde gerçekleştirilebilmesidir. Veritabanı işlemleri pgAdmin arayüzü ile
yönetilmiştir.

Mevcut aşamada veri kaynağı olarak AFAD tarafından yayımlanan güncel deprem kayıtları
kullanılmıştır. İlerleyen aşamalarda Kandilli Rasathanesi verilerinin de sisteme dahil
edilmesi planlanmaktadır.

\section{Sistem Akışı}

Projede oluşturulan sistem akışı temel olarak dört aşamadan oluşmaktadır:

\begin{enumerate}
    \item AFAD kaynağından deprem verilerinin alınması,
    \item alınan verilerin temizlenmesi ve ortak veri şemasına dönüştürülmesi,
    \item temizlenen verilerin PostgreSQL veritabanına aktarılması,
    \item veritabanı üzerinde analiz ve risk puanlama işlemlerinin gerçekleştirilmesi.
\end{enumerate}

İlk aşamada AFAD internet sitesindeki son depremler tablosu Python kullanılarak
okunmuştur. Ardından tarih-saat, enlem, boylam, derinlik ve büyüklük gibi alanlar
uygun veri tiplerine dönüştürülmüştür. Yer bilgisi alanı üzerinde il ve ilçe bilgisi
ayrıştırılmış, eksik veya hatalı kayıtlar temizlenmiştir. Bu işlem sonucunda elde
edilen temiz veri CSV biçiminde kaydedilmiş ve daha sonra PostgreSQL veritabanına
aktarılmıştır.

Veritabanına aktarılan kayıtlar üzerinde şehir bazında deprem sayısı, ortalama
büyüklük, maksimum büyüklük ve risk skoru hesaplamaları yapılmıştır. Böylece veri
toplama ile ilk karar destek çıktıları arasında izlenebilir ve tekrarlanabilir bir
iş akışı oluşturulmuştur.

\section{Belirlenen Veri Alanları}

Ara rapor aşamasında sistemde tutulacak temel veri alanları belirlenmiş ve uygulamada
kullanılmaya başlanmıştır. Bu alanlar hem olay bazlı kayıt takibi hem de bölgesel
özet analizler için yeterli bir temel oluşturmaktadır.

\begin{table}[H]
    \centering
    \caption{Deprem kayıtları için belirlenen temel veri alanları}
    \begin{tabular}{|l|p{8cm}|}
        \hline
        \textbf{Veri Alanı} & \textbf{Açıklama} \\
        \hline
        Kaynak ID & Verinin geldiği kurum bilgisi \\
        \hline
        Tarih & Deprem olayının gerçekleştiği tarih \\
        \hline
        Saat & Deprem olayının gerçekleştiği saat \\
        \hline
        Enlem & Depremin enlem değeri \\
        \hline
        Boylam & Depremin boylam değeri \\
        \hline
        Derinlik & Deprem odağının kilometre cinsinden derinliği \\
        \hline
        Büyüklük & Depremin sayısal büyüklük değeri \\
        \hline
        Yer Bilgisi & Deprem kaydındaki ham konum açıklaması \\
        \hline
        İl & Ayrıştırılmış il bilgisi \\
        \hline
        İlçe & Ayrıştırılmış ilçe bilgisi \\
        \hline
        Kayıt Kimliği & Veritabanındaki tekil kayıt kimliği \\
        \hline
    \end{tabular}
    \label{tab:verialanlari}
\end{table}

Bu alanlar sayesinde hem tekil deprem olaylarının saklanması hem de şehir bazında
özet sonuçlar üretilmesi mümkün hâle gelmiştir. İlerleyen aşamalarda veri kaynağı
çeşitlendiğinde büyüklük türü, kayıt kimliği ve ek türetilmiş göstergeler de daha
geniş biçimde kullanılabilecektir.

\section{Risk Analizi Yaklaşımı}

Projede deprem riski, tek bir değişkene bağlı olarak değil; birden fazla göstergenin
birlikte değerlendirilmesiyle yorumlanmaktadır. Ara rapor aşamasında uygulanan ilk
yaklaşımda şehir bazında toplam deprem sayısı, ortalama büyüklük ve maksimum büyüklük
değerleri kullanılmıştır. Bu göstergeler kullanılarak örnek bir risk skoru üretilmiştir.

Mevcut aşamada kullanılan örnek risk puanı aşağıdaki bileşenlere dayanmaktadır:

\begin{itemize}
    \item toplam deprem sayısı,
    \item ortalama büyüklük,
    \item maksimum büyüklük.
\end{itemize}

Uygulanan örnek hesaplama yaklaşımı şu şekildedir:

\[
Risk\ Skoru = (Toplam\ Deprem\ Sayısı \times 0.4) + (Ortalama\ Büyüklük \times 0.4) + (Maksimum\ Büyüklük \times 0.2)
\]

Bu formül nihai risk modeli değildir; ancak mevcut veri üzerinde bölgesel karşılaştırma
yapılabilmesi için başlangıç niteliğinde bir karar destek göstergesi sunmaktadır.
İlerleyen aşamalarda derinlik profili, eşik üstü deprem sayısı, son dönem aktivite
yoğunluğu ve gerekirse kümelenme temelli göstergeler de modele dahil edilecektir.

\chapter{Uygulama Durumu}

Ara rapor tarihine kadar gerçekleştirilen çalışmalar aşağıda özetlenmiştir:

\begin{itemize}
    \item Proje konusu netleştirilmiş ve problem tanımı oluşturulmuştur.
    \item AFAD ve Kandilli veri kaynakları incelenmiş, sistemde kullanılabilecek temel veri alanları belirlenmiştir.
    \item PostgreSQL tabanlı veritabanı yapısı oluşturulmuş; \texttt{data\_sources}, \texttt{earthquakes} ve \texttt{risk\_analysis} tabloları tasarlanmış ve oluşturulmuştur.
    \item AFAD internet sitesinden güncel deprem verilerini çekmek için Python tabanlı veri toplama modülü geliştirilmiştir.
    \item Çekilen veriler üzerinde tarih-saat dönüştürme, sayısal alan temizleme ve yer bilgisinden il/ilçe ayrıştırma işlemleri gerçekleştirilmiştir.
    \item Temizlenen veriler CSV formatına dönüştürülmüş ve PostgreSQL veritabanına aktarılmıştır.
    \item Veritabanına aktarılan kayıtlar üzerinde şehir bazında deprem sayısı, ortalama büyüklük ve maksimum büyüklük gibi ilk analiz sorguları çalıştırılmıştır.
    \item Bölgesel karşılaştırma amacıyla örnek bir risk puanlama yaklaşımı uygulanmış ve sonuçlar \texttt{risk\_analysis} tablosuna kaydedilmiştir.
    \item Elde edilen risk skorları Python ve \texttt{matplotlib} kullanılarak grafiksel olarak görselleştirilmiştir.
\end{itemize}

Mevcut durumda sistemin temel veri işleme hattı çalışır durumdadır. AFAD kaynağından alınan veriler okunabilmekte, temizlenebilmekte, veritabanına aktarılabilmekte ve temel analiz sorguları ile işlenebilmektedir. Bu durum, projenin yalnızca kavramsal tasarım aşamasında kalmadığını, aynı zamanda çalışan bir prototip düzeyine ulaştığını göstermektedir.

Bununla birlikte, proje henüz tamamlanmış değildir. Sonraki geliştirme adımları arasında aşağıdaki çalışmalar yer almaktadır:

\begin{itemize}
    \item Kandilli verilerinin sisteme dahil edilmesi,
    \item farklı veri kaynaklarının ortak veri şemasında birleştirilmesi,
    \item risk puanlama modelinin geliştirilmesi,
    \item zaman serisi ve kümelenme analizlerinin eklenmesi,
    \item web arayüzünün ilk prototipinin hazırlanması.
\end{itemize}

Gelinen noktada proje, başlangıç düzeyindeki planlama aşamasını aşmış; veri toplama, veri temizleme, veritabanına aktarma ve ilk analiz çıktılarının üretilmesi aşamasına ulaşmıştır.

\chapter{Sonuçlar ve Tartışma}

Ara rapor aşamasında sistemin nihai çıktıları henüz elde edilmemiştir. Bununla birlikte,
uygulama sürecinde geliştirilen veri toplama, temizleme ve analiz bileşenleri sayesinde
ilk somut sonuçlar üretilmiştir. AFAD kaynağından alınan deprem verileri PostgreSQL
veritabanına aktarılmış ve şehir bazında temel özet istatistikler elde edilmiştir.

Yapılan ilk analizlerde bazı şehirlerin veri kümesinde daha yoğun biçimde öne çıktığı
gözlenmiştir. Özellikle Kütahya, Muğla, Balıkesir, Kahramanmaraş ve Bolu illeri,
mevcut veri üzerinde hesaplanan örnek risk skorlarında üst sıralarda yer almıştır.
Bu sonuçlar, kullanılan risk formülünün toplam deprem sayısı, ortalama büyüklük ve
maksimum büyüklük bileşenlerini birlikte değerlendirmesine dayanmaktadır.
Risk skorlarının ilk 10 il için dağılımı Şekil~\ref{fig:risk_scores}'de gösterilmiştir.

\begin{figure}[H]
    \centering
    \includegraphics[width=0.88\textwidth,trim=2cm 1.5cm 2cm 1.5cm,clip]{risk_scores.png}
    \caption{İllere Göre Risk Skorları (İlk 10)}
    \label{fig:risk_scores}
\end{figure}

Elde edilen örnek sonuçlar, sistemin yalnızca deprem kayıtlarını depolayan bir yapı
olmadığını; aynı zamanda verileri işleyerek bölgesel karşılaştırma yapabilen bir karar
destek yaklaşımına dönüştürülebileceğini göstermektedir. Bu durum, projenin temel
hedefiyle uyumludur. Ancak bu aşamada kullanılan risk skoru, nihai ve bilimsel olarak
tamamlanmış bir deprem tehlike modeli olarak değerlendirilmemelidir. Mevcut yaklaşım,
projenin erken aşamasında analiz mantığını test etmek ve farklı bölgeler arasında ilk
karşılaştırmaları yapabilmek amacıyla oluşturulmuş başlangıç niteliğinde bir göstergedir.

Uygulama sürecinde bazı sınırlılıklar da gözlenmiştir. Öncelikle veri kaynağındaki alan
isimleri ve veri sunum biçimi zaman içerisinde değişebileceğinden, veri çekme ve temizleme
modüllerinin dayanıklı biçimde geliştirilmesi gerekmektedir. Ayrıca şu ana kadar yalnızca
AFAD verileri sisteme dahil edilmiştir; Kandilli verilerinin eklenmesi ile veri kapsamı
ve karşılaştırma gücü artırılacaktır. Bunun yanında mevcut risk puanlama yaklaşımına
derinlik profili, eşik üstü deprem sayısı, son dönem aktivite yoğunluğu ve kümelenme
analizi gibi ek bileşenlerin dahil edilmesi planlanmaktadır.

Sonuç olarak mevcut bulgular, projenin teknik açıdan uygulanabilir olduğunu ve resmi
deprem verilerinin tek sistem altında analiz edilmesinin anlamlı çıktılar üretebildiğini
göstermektedir. Bundan sonraki aşamada sistemin analitik kapasitesi geliştirilecek ve
elde edilen sonuçların daha güçlü görselleştirmelerle kullanıcıya sunulması sağlanacaktır.

\chapter{Sonuç}

Bu ara raporda, Türkiye deprem verilerine dayalı riskli bölge ve olay deseni analizi
yapacak bir karar destek sistemi için şimdiye kadar tamamlanan çalışmalar sunulmuştur.
Gelinen noktada proje problemi tanımlanmış, veri kaynakları değerlendirilmiş, sistemde
kullanılacak temel veri alanları belirlenmiş ve teknik altyapı oluşturulmuştur.

Ara rapor aşamasında PostgreSQL tabanlı veritabanı yapısı kurulmuş, AFAD verilerinin
Python kullanılarak alınması sağlanmış, veriler temizlenmiş ve veritabanına aktarılmıştır.
Ayrıca şehir bazında temel analiz sorguları uygulanmış ve örnek bir risk puanlama
yaklaşımı geliştirilmiştir. Elde edilen risk skorlarının grafiksel olarak gösterilmesiyle
birlikte, projenin yalnızca kavramsal düzeyde kalmadığı, çalışan bir prototip düzeyine
ulaştığı görülmüştür.

Mevcut değerlendirmeler, projenin hem teknik hem de akademik açıdan uygulanabilir
olduğunu göstermektedir. Türkiye'ye ait resmi deprem verilerinin tek sistem altında
toplanması ve analiz edilmesi, kullanıcıların deprem hareketliliğini daha bilinçli
yorumlamasına katkı sağlayabilecek bir yazılım ortaya çıkarma potansiyeline sahiptir.

Sonraki aşamada Kandilli verilerinin sisteme dahil edilmesi, veri kaynaklarının ortak
bir yapıda birleştirilmesi, risk puanlama modelinin geliştirilmesi, zaman serisi ve
kümelenme analizlerinin eklenmesi ve web tabanlı kullanıcı arayüzünün oluşturulması
hedeflenmektedir. Böylece proje, mevcut çalışan analiz altyapısından daha kapsamlı ve
kullanıcıya doğrudan karar desteği sunan bir sisteme dönüştürülecektir.

%-------------------- KAYNAKÇA --------------------
\begin{thebibliography}{99}

\bibitem{afad_harita}
Afet ve Acil Durum Yönetimi Başkanlığı,
\textit{Türkiye Deprem Tehlike Haritası},
Erişim Tarihi: 14.04.2026.

\bibitem{afad_son}
Afet ve Acil Durum Yönetimi Başkanlığı,
\textit{Son Depremler},
Erişim Tarihi: 14.04.2026.

\bibitem{kandilli_son}
Boğaziçi Üniversitesi Kandilli Rasathanesi ve Deprem Araştırma Enstitüsü,
\textit{Son Depremler},
Erişim Tarihi: 14.04.2026.

\bibitem{kandilli_yillik}
Boğaziçi Üniversitesi Kandilli Rasathanesi ve Deprem Araştırma Enstitüsü,
\textit{Yıllık Deprem Harita, Grafik ve Tabloları},
Erişim Tarihi: 14.04.2026.

\end{thebibliography}

\end{document}